\documentclass[11pt]{article}
\usepackage[margin=1.15in]{geometry}
\usepackage{amsmath,amssymb,amsthm,mathtools}
\usepackage{xcolor}
\usepackage[colorlinks=true,linkcolor=blue!60!black,citecolor=blue!60!black,urlcolor=blue!60!black]{hyperref}

\newtheorem{theorem}{Theorem}[section]
\newtheorem{proposition}[theorem]{Proposition}
\newtheorem{lemma}[theorem]{Lemma}
\newtheorem{corollary}[theorem]{Corollary}
\theoremstyle{definition}
\newtheorem{definition}[theorem]{Definition}
\theoremstyle{remark}
\newtheorem{remark}[theorem]{Remark}

\newcommand{\C}{\mathbb{C}}
\newcommand{\R}{\mathbb{R}}
\newcommand{\Z}{\mathbb{Z}}
\newcommand{\w}{\omega}
\newcommand{\Sp}{\mathrm{Sp}}
\newcommand{\dg}{\dagger}
\newcommand{\dd}{\ddagger}
\newcommand{\WF}{\operatorname{WF}}

\title{Gauge Reduction and the Completion Problem in Fourth-Order Gravity:\\
PU Pairing, Covariant Real Forms, and the Conformal Jordan Boundary}
\author{}
\date{}

\begin{document}
\maketitle

\begin{abstract}
We classify the quantizations of free scalar-free Einstein--Weyl gravity,
$\mathcal L=\sqrt{-g}\,(c_1R+\alpha R_{\mu\nu}^2+\beta R^2)$ with
$\alpha=-3\beta$, around flat space.  The algebra of linear
gauge-invariant observables carries a unique Poincar\'e-covariant complex
spectral quasifree functional; among translation-invariant, quasifree,
mode-local completions compatible with the reduced symplectic form, the
spectral condition, and massive-spin-$2$ covariance, there are exactly two
equivalence classes in the split phase: a positive pseudo-Hermitian
completion in which the full massive spin-$2$ multiplet has uniformly
rotated reality, and a Krein completion with standard gravitational
reality, in which the massive spin-$2$ irrep carries a uniform negative
sign relative to the positive massless graviton sector.  No
helicity-hybrid completion is Lorentz covariant, the two classes have
identical complex gauge-invariant correlation functions, and only the
Krein class continues through the conformal Jordan boundary $c_1\to0$,
where the positive class terminates and the induced functional on the
gauge-invariant quotient has a regular distributional limit.  The
mechanism is a stratification: diffeomorphism reduction preserves the
Pais--Uhlenbeck (PU) pairing of the fourth-order dynamics exactly in the
helicity-$\pm2$ sector, while the massive helicity-$\pm1,0$ polarizations
survive as unpaired ghost oscillators subject to a completion trilemma ---
positivity, the spectral condition, and standard reality: any two, never
all three.  On the conformal locus the transverse-traceless sector becomes
a $\Box^2$ Jordan block per polarization, the vector modes become ordinary
massless ghosts, and the scalar mode enters the kernel of the presymplectic
form and is removed by the emergent Weyl gauge symmetry, giving the
conformal count of six.  All finite-dimensional, variational, symplectic,
tensor-algebraic, and modewise distributional identities are independently
machine-verified.
\end{abstract}

\section{Introduction}\label{sec:intro}

Fourth-order gravity has two well-developed responses to the Ostrogradsky
ghost.  In the PT-symmetric program of Bender and Mannheim
\cite{BM2008PRD,Mannheim2021}, a positive metric operator renders each
Pais--Uhlenbeck (PU) mode pair unitary on a Hilbert space; in the
Krein-space program, developed for the pure fourth-order scalar by Bateman
and Turok \cite{BT2026}, the state space is indefinite and probabilities
are controlled by ghost parity.  A companion series
\cite{companion1,companion2,companion3} analyzed the scalar case
exhaustively: the positive metric is canonically reconstructed and
variationally selected; particle content decouples from metric geometry;
and the two programs determine the \emph{same} complex spectral vacuum
functional, differing only in the real form --- the involution and
completion --- placed upon it.

Gravity poses a question the scalar theory cannot: the completion must be
chosen \emph{jointly} for all polarizations of a gauge theory, subject to
diffeomorphism reduction and Lorentz covariance.  This paper answers it
for linearized scalar-free Einstein--Weyl gravity.  The logic runs
\begin{equation*}
\text{reduction}\ \rightarrow\ \text{completion trilemma}\ \rightarrow\
\text{covariance classification}\ \rightarrow\ \text{spectral bridge}\
\rightarrow\ \text{gauge-invariant correlator}\ \rightarrow\
\text{conformal boundary},
\end{equation*}
and the results are:

\begin{enumerate}
\item \emph{Reduction stratifies} (Section~\ref{sec:reduction}).  At fixed
spatial momentum the physical phase space is
\begin{equation}\label{eq:strat}
\Gamma_{\mathrm{phys}}(k)\ \cong\
2\,\Gamma_{\mathrm{PU}}(0,M)\ \oplus\ 2\,\Gamma_-(M)\ \oplus\ \Gamma_-(M),
\qquad M^2=c_1/\alpha:
\end{equation}
the two transverse-traceless (TT) polarizations are exact PU pairs
(massless plus massive branch, in the perfect-square-plus-Einstein form
with $\gamma=\alpha/2$); the massive helicity-$\pm1$ and $0$ modes are
\emph{unpaired} second-order ghost oscillators, because their massless
partners lie entirely in the diffeomorphism orbit.  PU pairing survives
gauge reduction exactly in helicity $\pm2$.

\item \emph{Completion trilemma} (Section~\ref{sec:trilemma}).  An
unpaired ghost oscillator admits a quarter-turn positive completion
($q\mapsto iq$, $p\mapsto-ip$, $\eta=e^{-\pi D}$, $D=\tfrac12(qp+pq)$)
whose physical adjoint flips the sign of the original variables.  The
three desiderata --- positive norm, positive energy (spectral condition),
and standard gravitational reality --- can be satisfied two at a time,
never all three.  All complex-symplectic real-form changes taking $H_-$ to
$H_+$ form a single quarter-turn coset $T_0\cdot SO(2,\C)$.

\item \emph{Covariance classification} (Section~\ref{sec:covariance}).
The five massive polarizations form one spin-$2$ irrep of the massive
little group; by Schur, any invariant Hermitian form on it is a multiple
of the identity, so its signature is uniform: \emph{no helicity-hybrid
completion is covariant}.  The total quarter-turn generator is
$SO(3)$-invariant, so a covariant uniform-positive completion exists; and
the quarter-turn on the ghost normal mode of the TT PU pair is itself an
admissible positive diagonalizer, hence lies in the metric family of
\cite{companion1} and, by orbit constancy \cite{companion3}, defines the
same vacuum functional as the Bender--Mannheim metric: the TT and
lower-helicity structures assemble.

\item \emph{Spectral bridge} (Section~\ref{sec:kernel}).  The reduced
spectral kernel, with residues fixed by the reduced symplectic form,
reassembles --- modulo pure-gauge bi-distributions --- into the covariant
massive spin-$2$ projector with a single normalization $\mathcal N=4/c_1$
and uniform ghost sign; the massless shell couples only to TT.  The
complex kernel is identical in both real forms.  Together with item 3 this
closes the bridge at the gravitational level: \emph{the two covariant
completions induce the same vacuum functional on the gauge-invariant
observable algebra}.

\item \emph{Gauge-invariant correlator} (Section~\ref{sec:weyl}).  The
linearized Weyl correlator is gauge independent, and the Weyl map
annihilates every $p_\mu p_\nu/M^2$ term of the massive projector: the
curvature kernel is $M$-regular and covariant, with Hadamard wavefront
directions; the fourth-order logarithm of the potential kernel appears in
the curvature correlator as its differentiated descendant.

\item \emph{Conformal Jordan boundary} (Section~\ref{sec:conformal}).  As
$c_1\to0$ at fixed $\alpha$: the TT divided difference converges
distributionally to the $\Box^2$ Jordan kernel; the vector modes have a
smooth massless-ghost limit with fixed commutator normalization; the
scalar direction enters the kernel of the presymplectic form and becomes
Weyl-gauge (the transformation $\delta h_{\mu\nu}=2\sigma\eta_{\mu\nu}$
is a symmetry exactly at $c_1=0$), giving the conformal count $4+2+0=6$.
The limit is regular for the functional \emph{induced on the common
gauge-invariant algebra after quotienting the emergent Weyl-null
direction} --- the observable algebra changes rank at the boundary.  The
split rapidity of the TT normal-mode decomposition underlying the positive
completion diverges exactly, $e^{r}=(\sqrt{k^2+M^2}+k)^2/M^2\sim4k^2/M^2$:
the positive real form terminates, the Krein form and the spectral
functional continue.
\end{enumerate}

The classification is summarized by the master theorem of
Section~\ref{sec:master}.  Throughout, gauge (Faddeev--Popov/BRST) ghosts
--- which encode the diffeomorphism quotient and are handled here by
explicit symplectic reduction --- are sharply distinguished from the
higher-derivative spin-$2$ ghost branch, which concerns the completion of
the physical excitations.

\paragraph{Verification.}
All finite-dimensional, variational, symplectic, tensor-algebraic, and
modewise distributional identities in this paper are independently
machine-verified: the second variation and sector reduction (checks
G1--G7), the trilemma and real-form classification (G8), the covariance
computations (G9), the spectral kernel and covariant reassembly (G10), the
Weyl-correlator cancellations (G11), and the sectorwise conformal limits
including the exact divergence law of the split decomposition (G12), in
the verification repository of \cite{companion1}.  The
representation-theoretic and microlocal conclusions --- existence of the
Hilbert/Krein completions, Poincar\'e covariance as a representation
statement, and wavefront-set equality --- follow from the analytic
arguments given in the text and the companion papers, not from the
symbolic engine.

\paragraph{Conventions.}
Signature $(+,-,-,-)$; curvature conventions fixed so that the healthy
Einstein normalization is $c_1=-1$ (checked on the graviton kinetic term);
$\alpha=-3\beta$ (scalar-free tuning); $M^2=c_1/\alpha>0$ for $\alpha<0$;
modes $h_{\mu\nu}\sim A(t)\cos kz+B(t)\sin kz$ at spatial momentum $k$
along $z$; $\w_1=\sqrt{k^2+M^2}$, $\w_2=k$.  PU data follow
\cite{companion1,companion2}: a PU pair with parameter $\gamma$ has
Lagrangian $\tfrac{\gamma}{2}(\ddot z^2-(\w_1^2+\w_2^2)\dot
z^2+\w_1^2\w_2^2z^2)$ and commutator normalization $E'''(0)=1/\gamma$.

\section{Reduction: the helicity stratification}\label{sec:reduction}

The second variation of
$\mathcal L=\sqrt{-g}(c_1R+\alpha R_{\mu\nu}^2+\beta R^2)$ about flat
space is computed exactly, sector by sector under the rotation group about
$k$; diffeomorphism invariance of each sector Lagrangian is verified (the
gauge variation is a total derivative).

\begin{theorem}[TT sector: exact PU blocks]\label{thm:tt}
For each of the two TT polarizations the mode Lagrangian is
\begin{equation}\label{eq:LTT}
L_{\mathrm{TT}}
=\frac{\alpha}{4}\big(\ddot A+k^2A\big)^2
-\frac{c_1}{4}\big(\dot A^2-k^2A^2\big)+\frac{d}{dt}(\cdots),
\end{equation}
a perfect square plus Einstein term, with Euler--Lagrange equation
$(\partial_t^2+k^2)(\partial_t^2+k^2+M^2)A=0$, $M^2=c_1/\alpha$: an exact
PU block with mass pair $(M,0)$ and normalization $\gamma=\alpha/2$.  The
sign $\gamma<0$ (for $c_1=-1$, $\alpha<0$) is the perfect-square
convention of \cite{BT2026} and renders the \emph{massless} branch healthy
and the massive spin-$2$ branch the ghost, as required by the flat-space
spectrum of quadratic gravity \cite{Stelle}.
\end{theorem}

\begin{theorem}[Lower helicities: PU pairing broken by gauge]
\label{thm:lower}
(i) Vector sector: with the gauge-invariant $w=k\,h_{tx}+\partial_t
h_{xz}$ (one per transverse direction),
\begin{equation}
L_V=\frac{\alpha}{4}\big(\dot w^2-k^2w^2\big)-\frac{c_1}{4}w^2
+\frac{d}{dt}(\cdots):
\end{equation}
a single second-order massive mode, $\ddot w+(k^2+M^2)w=0$, with
ghost-signed kinetic normalization $\mu_V=\alpha/2$; the massless PU
partner is pure diffeomorphism gauge and is removed by the quotient.
(ii) Scalar sector: $h_{tt}$ is auxiliary; at $\alpha=-3\beta$ the reduced
dynamics is the single massive mode $\ddot p+(k^2+M^2)p=0$ with effective
kinetic normalization $\mu_S=3c_1/2$ (ghost-signed, $\beta$-independent);
for generic couplings the sector is fourth order with scalaron mass
$m_0^2=-c_1/(2(\alpha+3\beta))$, decoupling exactly at the scalar-free
tuning.  (iii) Mode count: $2\times2+2+1=7=2+5$, and the stratification
\eqref{eq:strat} holds: PU pairing survives reduction exactly in helicity
$\pm2$.
\end{theorem}

\begin{definition}[Invariant observable algebra]\label{def:Ainv}
$\mathfrak A_{\mathrm{inv}}$ denotes the algebra of linear gauge-invariant
observables of the linearized theory --- equivalently, the field algebra
of the reduced phase space \eqref{eq:strat} --- with, at $c_1=0$, the
additional quotient by the kernel of the presymplectic form (the emergent
Weyl-null direction of Section~\ref{sec:conformal}).  All uniqueness and
continuation statements in this paper are statements about states and
kernels on $\mathfrak A_{\mathrm{inv}}$; potential-field kernels
$W^+_h$ are determined only modulo pure-gauge bi-distributions and are
used as computational representatives.
\end{definition}

\begin{remark}[Polarization-enhanced stabilizer]
For the doubled TT block the normal-form stabilizer Lie algebra jumps from
$\dim_\R 4$ (one PU pair; $so(2,\C)^2$ \cite{companion1}) to $\dim_\R16$
(computed by null-space methods): frequency degeneracy across
polarizations enlarges the metric ambiguity.  $SO(2)$ helicity covariance
and Schur reduce the covariant TT metric to $S_+\otimes I_2$, the unique
helicity-covariant minimum-distortion positive PU diagonalizer
\cite{companion2}.
\end{remark}

\section{The completion trilemma}\label{sec:trilemma}

Each unpaired ghost oscillator, $H_-=-\tfrac12(p^2+\Omega^2q^2)$, poses
the completion problem in miniature.

\begin{theorem}[Unpaired-ghost completion]\label{thm:trilemma}
Let $D=\tfrac12(qp+pq)$ and $\rho=e^{-\frac{\pi}{2}D}$.  Then:
(i) $\rho q\rho^{-1}=iq$, $\rho p\rho^{-1}=-ip$, and
$\rho H_-\rho^{-1}=+\tfrac12(p^2+\Omega^2q^2)$;
$\eta=\rho^{\dg}\rho=e^{-\pi D}$ is formally positive.
(ii) The physical adjoint is $q^{\dd}=\eta^{-1}q^{\dg}\eta=-q$,
$p^{\dd}=-p$: the original amplitude is $\dd$-anti-Hermitian; $iq$ is the
observable.  (iii) Trilemma: positive norm, positive energy, and standard
reality can be realized two at a time, never all three ---
\begin{center}
\begin{tabular}{c|c|c|c}
completion & positive norm & positive energy & standard reality\\
\hline
ordinary Hilbert & \checkmark & $\times$ & \checkmark\\
Krein & $\times$ & \checkmark & \checkmark\\
quarter-turn (PT contour) & \checkmark & \checkmark & $\times$
\end{tabular}
\end{center}
(for the first row, $p^2+\Omega^2q^2$ built from a self-adjoint canonical
pair has positive unbounded spectrum, so $H_-$ is unbounded below).
(iv) Classification: all complex-symplectic real-form changes taking $H_-$
to $H_+$ form the single coset
$\{T\in\Sp(2,\C):TA_+T^{-1}=-A_+\}=T_0\cdot SO(2,\C)$,
$T_0=\operatorname{diag}(i,-i)$.
\end{theorem}

An unpaired gravitational ghost therefore cannot simultaneously satisfy
positivity, the spectral condition, and the standard real structure.  The
choice among the three completions is exactly a choice of real form.

\section{Covariance classification}\label{sec:covariance}

Sectorwise choices are constrained globally: the massive helicities
$(\pm2,\pm1,0)$ form one irreducible spin-$2$ multiplet of the massive
little group.

\begin{theorem}[No covariant hybrid]\label{thm:nohybrid}
The space of $SO(3)$-invariant Hermitian forms on the five-dimensional
spin-$2$ irrep is one-dimensional, $\eta_{\mathrm{mass}}=cI_5$ (computed
by solving the commutant equations).  Hence any covariant Hermitian form
has uniform signature, and the helicity-hybrid assignment
$(+,+,-,-,-)$ --- TT-positive with lower-helicity Krein --- violates
invariance explicitly and is not Lorentz covariant.
\end{theorem}

\begin{theorem}[Two covariant real forms]\label{thm:tworeal}
(i) The total quarter-turn generator
$D_{\mathrm{tot}}=\sum_{a=1}^{5}\tfrac12(q_ap_a+p_aq_a)$ over the massive
polarization oscillators is $SO(3)$-invariant; hence
$\eta_{\mathrm{mass}}=e^{-\pi D_{\mathrm{tot}}}$ defines a covariant
uniform-positive pseudo-Hermitian completion with uniformly rotated
massive reality (the massive multiplet is $\dd$-anti-Hermitian; $i$ times
the multiplet is observable).  (ii) The Krein completion with standard
gravitational reality is covariant, and by Theorem~\ref{thm:nohybrid} its
indefiniteness lives \emph{between} irreps, never inside the massive
multiplet: the invariant form on the massive spin-$2$ one-particle irrep
is definite, and the Krein structure is
\begin{equation}
\mathcal H_{\mathrm{1p}}
=\mathcal H_{0,+2}\ \oplus\ \mathcal H_{M,-5},
\qquad\text{signature } (+,+;\,-,-,-,-,-),
\end{equation}
the direct sum of the positive massless graviton sector and the massive
spin-$2$ irrep carrying uniform negative sign.  (iii) Assembly: the
quarter-turn applied to the ghost normal
mode of the TT PU pair is an admissible positive diagonalizer
($T'^{\top}JT'=J$ and $T'^{\top}G_{\mathrm{PU}}T'$ real positive
definite), hence lies in the solution family $S_+\cdot\mathrm{Stab}$ of
\cite{companion1}; by orbit constancy \cite{companion3} it defines the
same vacuum functional as the Bender--Mannheim metric.  The TT and
lower-helicity structures of the positive class are mutually consistent.
\end{theorem}

Stabilizer choices (the enlarged TT stabilizer of
Section~\ref{sec:reduction}, the $SO(2,\C)$ freedom of
Theorem~\ref{thm:trilemma}(iv), the $W$-freedom of \cite{companion1})
produce representatives \emph{within} these classes, not additional
physical classes: by orbit constancy the vacuum ray, and hence the
quasifree functional, is unchanged.

\begin{remark}[Field-level meaning of the positive class]\label{rem:field}
At field level the positive completion is constructed in the order:
(1) modewise complex-symplectic real-form change
(Theorem~\ref{thm:trilemma}); (2) positive one-particle Hermitian
structure on the polarization spaces (Theorems~\ref{thm:nohybrid},
\ref{thm:tworeal}); (3) the corresponding quasifree (GNS)
representation.  The metric-operator notation
$\eta=e^{-\pi D_{\mathrm{tot}}}$ is formal shorthand valid on a common
algebraic domain; it is \emph{not} asserted to be a globally defined
invertible operator on the na\"ive Fock space --- indeed, by the
disjointness theorems of \cite{companion2,companion3}, the positive
quasifree representation of the field theory is disjoint from the na\"ive
Fock representation, and the representation is defined by step (3), not by
conjugating inside a pre-existing Hilbert space.
\end{remark}

\section{The reduced spectral kernel and the bridge}\label{sec:kernel}

\begin{theorem}[Helicity kernel with symplectic residues]\label{thm:kernel}
The reduced spectral (positive-frequency) kernel is
\begin{align}
W^+_{\mathrm{TT}}(t)
&=\frac{1}{\gamma M^2}
\left[\frac{e^{-i\w_1t}}{2\w_1}-\frac{e^{-i\w_2t}}{2\w_2}\right],
\qquad \gamma=\frac{\alpha}{2},\\
W^+_{V}(t)=\frac{e^{-i\Omega t}}{2\mu_V\Omega},
\quad
W^+_{S}(t)&=\frac{e^{-i\Omega t}}{2\mu_S\Omega},
\qquad
\Omega=\sqrt{k^2+M^2},\quad \mu_V=\frac{\alpha}{2},\ \mu_S=\frac{3c_1}{2},
\end{align}
each a bisolution with the commutator fixed by the reduced symplectic
form ($E'''(0)=1/\gamma$ for TT; $E'(0)=-1/\mu$ for the second-order
modes; the residues are \emph{derived}, not assumed from propagator
guesses).
\end{theorem}

\begin{theorem}[Covariant reassembly]\label{thm:reassembly}
Let $\Pi^{(2,M)}_{\mu\nu\rho\sigma}
=\tfrac12(P_{\mu\rho}P_{\nu\sigma}+P_{\mu\sigma}P_{\nu\rho})
-\tfrac13P_{\mu\nu}P_{\rho\sigma}$, $P_{\mu\nu}=\eta_{\mu\nu}-p_\mu
p_\nu/M^2$, on the massive shell.  The gauge-invariant contractions are
\begin{equation}
\Pi_{xy,xy}=\tfrac12,\qquad
c\,\Pi\,\bar c\big|_{w}=\frac{(E^2-k^2)^2}{2M^2}=\frac{M^2}{2},\qquad
\Pi\big|_{(h_{xx}+h_{yy})/2}=\tfrac16 ,
\end{equation}
where the complex-basis vector invariant is
$\tilde w=-ik\,h_{tx}-iE\,h_{xz}$.  These match the reduced residue ratios
$1:M^2:\tfrac13$ of Theorem~\ref{thm:kernel} exactly, with the single
overall normalization $\mathcal N=4/c_1$ and uniform ghost sign: the five
massive residues assemble into one $SO(3)$-invariant projector.
Consequently the potential kernel, as an equivalence class
\begin{equation}
[\widetilde W^+_h]\ \in\
\mathcal D'(\mathrm{Sym}^2\otimes\mathrm{Sym}^2)\big/
\{\text{pure-gauge bi-distributions}\},
\end{equation}
has the exact representative
\begin{equation}\label{eq:exactkernel}
\widetilde W^+_{\mu\nu\rho\sigma}(p)
=\frac{\mathcal N}{M^2}\,\theta(p^0)\left[
\Pi^{(2,M)}_{\mu\nu\rho\sigma}(p)\,\delta(p^2-M^2)
-\Pi^{(2,0)}_{\mu\nu\rho\sigma}(p;n)\,\delta(p^2)
\right],
\qquad \mathcal N=\frac{4}{c_1},
\end{equation}
where $\Pi^{(2,0)}(p;n)$ is the massless spin-$2$ projector in any
reference frame $n$; the class $[\widetilde W^+_h]$ is
frame-independent because changing $n$ changes $\Pi^{(2,0)}$ by a
pure-gauge kernel.  The massless shell couples only to TT (the vector and
scalar invariants contract to zero with the massless polarizations), and
the $1/M^2$ prefactor is the PU divided-difference normalization
$1/(\gamma M^2)$ of Theorem~\ref{thm:kernel} in covariant form.
Equivalently and frame-freely: the exact curvature kernel
$\widetilde{\mathcal W}^+_{CC}=\mathcal D_C\widetilde W^+_h\mathcal
D_C^{\top}$ of Section~\ref{sec:weyl} contains no gauge terms and no
singular projector pieces, and determines the state on
$\mathfrak A_{\mathrm{inv}}$.
\end{theorem}

\begin{theorem}[Real-form independence; the gravitational bridge]
\label{thm:bridge4}
The complex spectral kernel is identical in the two covariant real forms:
for each unpaired ghost the quarter-turn completion's physical Wightman
function equals the Krein spectral value,
$\langle(i\tilde w)(t)\,(i\tilde w)(0)\rangle
=e^{-i\Omega t}/(2\mu\Omega)$, and for the TT pairs the statement is the
orbit constancy and spectral identification of
\cite{companion3}.  Hence the positive pseudo-Hermitian and Krein
completions of scalar-free Einstein--Weyl gravity induce the same complex
quasifree functional on the gauge-invariant observable algebra
$\mathfrak A_{\mathrm{inv}}$; they differ in involution and completion
only.
\end{theorem}

\section{The gauge-invariant Weyl correlator}\label{sec:weyl}

Curvature observables remove both gauge dependence and projector
singularities.

\begin{theorem}[Weyl correlator]\label{thm:weyl}
Let $C^{(1)}[h]$ be the linearized Weyl tensor (traceless; validated
symbolically) and $\mathcal
W^+=\mathcal D_C\,W^+_h\,\mathcal D_C^{\top}$ the induced curvature
kernel.  Then: (i) the linearized Riemann map annihilates pure gauge
$h=p\otimes\xi+\xi\otimes p$ identically, so $\mathcal W^+$ is gauge
independent; (ii) the full Weyl--Weyl contraction of $\Pi^{(2,M)}$ equals
that of $\Pi_0$ (all $P\to\eta$): every $p_\mu p_\nu/M^2$ and $/M^4$ term
of the massive projector is annihilated --- the curvature kernel is
$M$-regular, and the five massive helicities enter only through the
$SO(3)$-invariant flat projector; (iii) both real forms give the same
complex curvature functional (Theorem~\ref{thm:bridge4}); (iv)
$\WF(\mathcal W^+)=\mathcal C^+$, the standard positive-frequency Hadamard
set with tensor polarization constraints.  The fourth-order logarithmic
singularity belongs to the potential kernel, $W_h^+\sim\log\rho$, while
$\mathcal W^+_{CC}\sim\partial^4\log\rho$ is its differentiated
descendant.  Contact terms (momentum-space polynomials) are separated from
the nonlocal two-point distribution.
\end{theorem}

\begin{proof}[Proof of (iv)]
Since $\mathcal W^+=\mathcal D_C W^+_h\mathcal D_C^\top$ with $\mathcal
D_C$ a differential operator, $\WF(\mathcal W^+)\subseteq\WF(W^+_h)
=\mathcal C^+$ \cite{companion3}.  For the reverse inclusion, the
principal symbol of $\mathcal D_C$ at a covector
$(p;-p)\in\mathcal C^+$ is the quadratic-in-$p$ linearized-Weyl symbol
$\sigma_C(p)$, which is nonvanishing on the physical spin-$2$
polarization subspace over the characteristic set: on the massive shell
$\sigma_C(p)\Pi^{(2,M)}\sigma_C(p)^\top$ has the nonzero full contraction
of Theorem~\ref{thm:weyl}(ii), and on the massless shell $\sigma_C$ is
nonzero on TT polarizations (the linearized Weyl tensor of a TT wave does
not vanish).  Hence no direction of $\mathcal C^+$ carried by $W^+_h$ on
physical polarizations is annihilated at symbol level, and
$\WF(\mathcal W^+)=\mathcal C^+$.
\end{proof}

\section{The conformal Jordan boundary}\label{sec:conformal}

Take $c_1\to0$ at fixed $\alpha$ (so $M^2\to0$) \emph{before} imposing any
completion.

\begin{theorem}[Sectorwise conformal limit]\label{thm:conformal}
(i) TT: distributionally,
$[\delta(p^2-M^2)-\delta(p^2)]/M^2\to-\delta'(p^2)$; per polarization the
kernel converges to the $\Box^2$ Jordan kernel
$-(1+ikt)e^{-ikt}/(4k^3)$ (four TT configuration modes).
(ii) Vector: $\delta(p^2-M^2)\to\delta(p^2)$ smoothly with \emph{fixed}
commutator normalization $\mu_V=\alpha/2$: ordinary second-order massless
ghost modes, not Jordan partners.
(iii) Scalar: the kinetic normalization $3c_1/4\to0$, and the Weyl
transformation $\delta_\sigma h_{\mu\nu}=2\sigma\eta_{\mu\nu}$ is a gauge
symmetry of the scalar-free action exactly at $c_1=0$ (and not for
$c_1\neq0$): the scalar direction enters the kernel of the presymplectic
form and is removed by the emergent Weyl gauge symmetry.
The conformal count is $4+2+0=6$.
\end{theorem}

\begin{remark}[Regularity holds on the quotient, not before it]
\label{rem:quotient}
The raw scalar kernel $W^+_S=e^{-i\Omega t}/(2\mu_S\Omega)$ with
$\mu_S=3c_1/2$ diverges like $1/c_1$: the full reduced-coordinate kernel
of the split phase has \emph{no} regular limit.  What converges is the
family of presymplectic phase spaces
\begin{equation}
(\Gamma_{c_1},\Omega_{c_1})\ \longrightarrow\
(\Gamma_0/\ker\Omega_0,\ \bar\Omega_0),
\end{equation}
together with the induced functional on the common gauge-invariant algebra
(Definition~\ref{def:Ainv}): the divergent direction is exactly the one
that enters $\ker\Omega_0$ and is quotiented by the emergent Weyl
symmetry.  The observable algebra changes rank at the boundary; the
regular-limit statement of Theorem~\ref{thm:master} refers to the
functional on $\mathfrak A_{\mathrm{inv}}$ after this quotient.
\end{remark}

\begin{center}
\begin{tabular}{c|c|c|c}
 & $\pm2$ & $\pm1$ & $0$\\
\hline
$c_1\neq0$ & PU pair & massive ghost & massive ghost\\
$c_1=0$ & $\Box^2$ Jordan & massless ghost & Weyl gauge/null
\end{tabular}
\end{center}

\begin{theorem}[Termination of the positive real form]\label{thm:term}
The uniform positive completion is built on the split normal-mode
decomposition of the TT PU pairs, whose rapidity is
\begin{equation}
r(k,M)=\log\frac{(\w_1+\w_2)^2}{\w_1^2-\w_2^2}
=\log\frac{\big(\sqrt{k^2+M^2}+k\big)^2}{M^2}
=\log\frac{4k^2}{M^2}+O\!\left(\frac{M^2}{k^2}\right).
\end{equation}
As $M\to0$ at fixed $k$, $r\to\infty$ and the symplectic normal-mode
transformation degenerates exactly as
$\operatorname{cond}(N)=e^{r}\,(1+o(1))\sim4k^2/M^2$; at the boundary
the dynamical map is a Jordan block and, by the Jordan obstruction of
\cite{companion1}, no positive-definite invariant form exists.  The
positive real form therefore has no invariant positive continuation
through $c_1=0$; the spectral functional and the Krein real form continue
distributionally.  (Numerically, $\operatorname{cond}(N)
=7.6,\,47.5,\,403,\,4447$ at $M=1,\,0.3,\,0.1,\,0.03$, $k=1$, with
$\operatorname{cond}(N)/e^{r}\to1$; this serves as a regression check of
the exact law, not as the proof.)
\end{theorem}

\section{The master theorem}\label{sec:master}

\begin{theorem}[Classification]\label{thm:master}
Consider free scalar-free Einstein--Weyl gravity linearized about
Minkowski space.  Let $\mathfrak A_{\mathrm{inv}}$ be the algebra of
linear gauge-invariant observables, equivalently the reduced field
algebra, with the additional Weyl-null quotient imposed at $c_1=0$
(Definition~\ref{def:Ainv}).  Within the class of translation-invariant,
quasifree, mode-local constructions compatible with the reduced symplectic
form, the spectral condition, and the Poincar\'e action on the physical
polarization spaces, the following hold.

For $c_1\neq0$, the reduced complex theory determines a unique
Poincar\'e-covariant spectral quasifree functional on
$\mathfrak A_{\mathrm{inv}}$.  It admits exactly two inequivalent
covariant real-form completions, modulo symplectic stabilizers commuting
with the dynamics:
\begin{enumerate}
\item a positive pseudo-Hermitian completion in which the full massive
spin-$2$ multiplet has uniformly rotated reality; and
\item a Krein completion in which the gravitational field retains its
standard reality and the massive spin-$2$ one-particle irrep carries
uniform negative sign relative to the positive massless graviton sector.
\end{enumerate}
No completion assigning different signatures or real structures to
different helicities of the massive spin-$2$ irrep is Poincar\'e
covariant.  The two admissible real forms induce the same complex
quasifree functional, and hence the same complex correlation functions on
$\mathfrak A_{\mathrm{inv}}$; they differ only in involution and
completion.

As $c_1\to0$, the induced functional on the gauge-invariant quotient has
a regular distributional limit.  The transverse-traceless sector becomes a
pair of $\Box^2$ Jordan systems, the vector sector becomes two ordinary
second-order massless ghost modes, and the scalar sector enters the kernel
of the presymplectic form and is removed by the emergent Weyl gauge
symmetry.  The positive real form has no invariant positive-definite
continuation through this boundary, whereas the Krein real form continues
to the resulting Jordan theory.
\end{theorem}

\begin{proof}
Uniqueness of the functional on $\mathfrak A_{\mathrm{inv}}$:
Theorems~\ref{thm:kernel}, \ref{thm:reassembly} (residues forced by the
reduced symplectic form and covariance; the potential representative is
fixed exactly modulo pure-gauge kernels by \eqref{eq:exactkernel}).  Two
real forms and no hybrid:
Theorems~\ref{thm:trilemma}--\ref{thm:tworeal} and \ref{thm:nohybrid}.
Identity of complex correlators: Theorems~\ref{thm:bridge4},
\ref{thm:weyl}.  Stabilizer statement: orbit constancy
\cite{companion3} with the enlarged stabilizers of
Section~\ref{sec:reduction}.  Boundary: Theorems~\ref{thm:conformal},
\ref{thm:term} and Remark~\ref{rem:quotient}.
\end{proof}

\section{Discussion}\label{sec:discussion}

\paragraph{Relation to the two programs.}
The theorem places the positive-PT quantization \cite{BM2008PRD,
Mannheim2021} and the Krein quantization on the two sides of
a mathematically controlled classification: they are the two covariant
real forms of one reduced complex spectral theory, agreeing on every
complex gauge-invariant correlation function, and distinguished by which
member of the trilemma they sacrifice --- standard gravitational reality
(positive form) or positivity of the norm (Krein form).  A terminological
caution: Bateman and Turok \cite{BT2026} quantize the scalar
perfect-square theory, not this gravitational field; the second class is
therefore a \emph{Bateman--Turok-type Krein real form} --- the
gravitational extension of their construction --- rather than their
construction itself.  The conformal locus is one-sided: only the Krein
class crosses it.  In this precise sense Mannheim's construction is a
split-phase real form and the Bateman--Turok-type Krein form is its
resonant continuation.

\paragraph{Boundary truncation is a third, distinct construction.}
Maldacena's boundary-condition selection of the Einstein branch from
conformal gravity \cite{Maldacena} removes one branch of the solution
algebra, $\ker(L_1L_2)\to\ker L_1$; the completions classified here retain
the full fourth-order solution algebra and change its inner-product
structure.  The two mechanisms answer the ghost problem differently and
should not be identified; on Einstein backgrounds with $\Lambda\neq0$ the
critical locus shifts \cite{LuPope} and pure conformal gravity contains a
partially massless sector \cite{DeserWaldron}, so the flat-space phase
diagram above is the $\Lambda\to0$ slice of a richer picture we do not
analyze here.

\paragraph{Interaction diagnostic (outlook).}
The free classification fixes the arena for the interacting question: for
the positive real form, whether the rotated massive reality is compatible
with the cubic Weyl vertices; for the Krein form, whether a gravitational
ghost parity commutes with the interaction and with BRST symmetry, as its
scalar analogue does in the perfect-square theory \cite{BT2026}.  Both are
sharp, computable questions --- the natural first test is
$[\mathcal P_{\mathrm{ghost}},S_{\mathrm{int}}]=0$ at cubic order --- and
they decide whether either real form supports an interacting unitarity
statement.  We have not performed these computations; the free theorem is
neutral on the outcome.

\begin{thebibliography}{12}

\bibitem{Stelle}
K.~S.~Stelle,
Phys.\ Rev.\ D \textbf{16} (1977) 953.

\bibitem{BM2008PRD}
C.~M.~Bender and P.~D.~Mannheim,
Phys.\ Rev.\ D \textbf{78} (2008) 025022.

\bibitem{Mannheim2021}
P.~D.~Mannheim,
\emph{Solution to the ghost problem in higher-derivative gravity},
arXiv:2109.12743.

\bibitem{BT2026}
S.~Bateman and N.~Turok,
\emph{Escape from Ostrogradsky via hidden ghost parity},
arXiv:2607.00096.

\bibitem{companion1}
\emph{Canonical positive symplectic diagonalization of the Pais--Uhlenbeck
oscillator}, companion paper and verification repository, 2026.

\bibitem{companion2}
\emph{The Pais--Uhlenbeck metric as a minimum-distortion principle, and
the representation problem for the fourth-order field}, companion paper,
2026.

\bibitem{companion3}
\emph{The universal vacuum of the fourth-order scalar field: metric
orbits, Fock sectors, and the Krein boundary}, companion paper, 2026.

\bibitem{LuPope}
H.~L\"u and C.~N.~Pope,
\emph{Critical gravity in four dimensions},
Phys.\ Rev.\ Lett.\ \textbf{106} (2011) 181302, arXiv:1101.1971.

\bibitem{Maldacena}
J.~Maldacena,
\emph{Einstein gravity from conformal gravity},
arXiv:1105.5632.

\bibitem{DeserWaldron}
S.~Deser, E.~Joung and A.~Waldron,
\emph{Partial masslessness and conformal gravity},
J.\ Phys.\ A \textbf{46} (2013) 214019, arXiv:1208.1307.

\bibitem{Holdom}
B.~Holdom,
Nucl.\ Phys.\ B \textbf{1000} (2024) 116696.

\bibitem{Riegert}
R.~J.~Riegert,
Phys.\ Lett.\ B \textbf{134} (1984) 56.

\end{thebibliography}

\end{document}
